% Options for packages loaded elsewhere
\PassOptionsToPackage{unicode}{hyperref}
\PassOptionsToPackage{hyphens}{url}
\documentclass[
  ignorenonframetext,
]{beamer}
\newif\ifbibliography
\usepackage{pgfpages}
\setbeamertemplate{caption}[numbered]
\setbeamertemplate{caption label separator}{: }
\setbeamercolor{caption name}{fg=normal text.fg}
\beamertemplatenavigationsymbolsempty
% remove section numbering
\setbeamertemplate{part page}{
  \centering
  \begin{beamercolorbox}[sep=16pt,center]{part title}
    \usebeamerfont{part title}\insertpart\par
  \end{beamercolorbox}
}
\setbeamertemplate{section page}{
  \centering
  \begin{beamercolorbox}[sep=12pt,center]{section title}
    \usebeamerfont{section title}\insertsection\par
  \end{beamercolorbox}
}
\setbeamertemplate{subsection page}{
  \centering
  \begin{beamercolorbox}[sep=8pt,center]{subsection title}
    \usebeamerfont{subsection title}\insertsubsection\par
  \end{beamercolorbox}
}
% Prevent slide breaks in the middle of a paragraph
\widowpenalties 1 10000
\raggedbottom
\AtBeginPart{
  \frame{\partpage}
}
\AtBeginSection{
  \ifbibliography
  \else
    \frame{\sectionpage}
  \fi
}
\AtBeginSubsection{
  \frame{\subsectionpage}
}
\usepackage{iftex}
\ifPDFTeX
  \usepackage[T1]{fontenc}
  \usepackage[utf8]{inputenc}
  \usepackage{textcomp} % provide euro and other symbols
\else % if luatex or xetex
  \usepackage{unicode-math} % this also loads fontspec
  \defaultfontfeatures{Scale=MatchLowercase}
  \defaultfontfeatures[\rmfamily]{Ligatures=TeX,Scale=1}
\fi
\usepackage{lmodern}
\ifPDFTeX\else
  % xetex/luatex font selection
\fi
% Use upquote if available, for straight quotes in verbatim environments
\IfFileExists{upquote.sty}{\usepackage{upquote}}{}
\IfFileExists{microtype.sty}{% use microtype if available
  \usepackage[]{microtype}
  \UseMicrotypeSet[protrusion]{basicmath} % disable protrusion for tt fonts
}{}
\makeatletter
\@ifundefined{KOMAClassName}{% if non-KOMA class
  \IfFileExists{parskip.sty}{%
    \usepackage{parskip}
  }{% else
    \setlength{\parindent}{0pt}
    \setlength{\parskip}{6pt plus 2pt minus 1pt}}
}{% if KOMA class
  \KOMAoptions{parskip=half}}
\makeatother
\setlength{\emergencystretch}{3em} % prevent overfull lines
\providecommand{\tightlist}{%
  \setlength{\itemsep}{0pt}\setlength{\parskip}{0pt}}
% Auto-generated by infrastructure/rendering/_pdf_latex_helpers.py::extract_math_font_preamble.
% Minimal header passed to Pandoc via -H so Beamer slides pick up the same math font
% as the combined-PDF path. Adjust the manuscript's preamble.md to change the font globally.
\usepackage{unicode-math}
\setmathfont{latinmodern-math.otf}
\usepackage{bookmark}
\IfFileExists{xurl.sty}{\usepackage{xurl}}{} % add URL line breaks if available
\urlstyle{same}
\hypersetup{
  hidelinks,
  pdfcreator={LaTeX via pandoc}}

\author{\texorpdfstring{}{}}
\date{}

\begin{document}

\section{Conclusion: Elevating Language Models from Vectors to Enriched
Category Frameworks}\label{sec:conclusion}

\begin{frame}[fragile,allowframebreaks]{What This Paper Actually Did:
Eleven Concrete Deliverables}
\protect\phantomsection\label{what-this-paper-actually-did-eleven-concrete-deliverables}
This paper developed a unified category-theoretic framework for
linguistic case systems, synthesizing \textbf{five formal
traditions}---typological, type-logical, distributional,
enriched-categorical, and topos-theoretic---together with a
\textbf{sixth strand}, the biolinguistic and neurocomputational
interface (ROSE and related oscillatory accounts;
\autoref{sec:introduction}, \autoref{sec:cognitive-integration}), and
embedding the whole within an active inference model of cognition. Our
eleven principal contributions are:

\begin{quote}
\textbf{Notation}: Each entry is labelled \textbf{C}\_n\_ (\emph{n}th
contribution, \textbf{C1}--\textbf{C11}) or \textbf{F}\_n\_ (\emph{n}th
future direction, \textbf{F1}--\textbf{F8}).
\end{quote}

\begin{block}{C1: Case Categories as a Formal Algebraic Framework}
\protect\phantomsection\label{c1-case-categories-as-a-formal-algebraic-framework}
The review formalized case systems as categories with case roles as
objects and grammatical relations as morphisms
(\autoref{sec:case-systems}). This formalization captures the full
typological range of alignment systems---nominative-accusative,
ergative-absolutive, active-stative, tripartite, and fluid-S---within a
single algebraic framework, with alignment functors providing
structure-preserving mappings between systems. By tracing the lineage
from Fillmore's {[}-@fillmore1968case{]} deep cases through Jakobson's
binary distinctive features and Dowty's {[}-@dowty1991thematic{]}
proto-roles, the analysis demonstrated how enriched (weighted) morphisms
link the categorical formalization to the gradient nature of thematic
role assignment.
\end{block}

\begin{block}{C2: String Diagrams for Case Derivation Visualization}
\protect\phantomsection\label{c2-string-diagrams-for-case-derivation-visualization}
Building on Joyal and Street's {[}-@joyalstreet1991geometry{]} string
diagram formalism and its application in categorial grammar
(\autoref{sec:categorial-grammar}), we showed how case-marked noun
phrases receive type-logical assignments that are fully visualizable as
string diagrams. The Curry--Howard correspondence ensures that syntactic
well-formedness guarantees semantic compositionality, and the
diagrammatic format provides Shimojima's {[}-@shimojima1996reasoning{]}
``free ride'' inferences---conclusions about argument structure that are
perceptually available from the diagram without explicit computation. We
further demonstrated that passivization reduces to a \emph{type
permutation} (a Swap operation in the pregroup category), making voice
alternation visible as a topological feature of the string diagram.
\end{block}

\begin{block}{C3: Case-Marked DisCoCat, the Distributional--Formal
Synthesis, and Discourse Extension}
\protect\phantomsection\label{c3-case-marked-discocat-the-distributionalformal-synthesis-and-discourse-extension}
We extended the DisCoCat framework (\autoref{sec:categorical-semantics})
with case-typed noun spaces and alignment-sensitive meaning functors,
and showed how the recent DisCoCirc {[}@defelice2022discocirc{]} and
QNLP {[}@lorenz2021lambeq{]} developments extend this analysis to
discourse-level structure and quantum hardware respectively. We
formalized the \emph{compact closure axiom} (snake equation) that
underpins pregroup reductions, demonstrating that the cup-cap zigzag
identity provides a visual coherence proof with genuine cognitive
significance. Diagram complexity metrics---normal form and
depth---provide a quantitative bridge between the type-logical and
distributional perspectives on linguistic structure. A key contribution
is the demonstration that DisCoCat constitutes the \emph{algebraic
formalization} of the distributional programme that modern large
language models---from Word2Vec {[}@mikolov2013efficient{]} through BERT
{[}@devlin2019bert{]} and GPT {[}@radford2018improving{]}---implement
empirically: the categorical meaning functor is the principled version
of the composition that transformer attention mechanisms learn from data
{[}@vaswani2017attention{]}. Case categories can serve as the structural
backbone of compositional models of meaning at all levels---word,
sentence, discourse, and dialogue.
\end{block}

\begin{block}{C4: Enriched Cases, Categorical Magnitude, and Information
Theory}
\protect\phantomsection\label{c4-enriched-cases-categorical-magnitude-and-information-theory}
Through Bradley et al.'s {[}-@fritz2021enriched{]} enrichment framework
and Bradley's {[}-@bradley2021entropy{]} information-theoretic analysis
(\autoref{sec:enriched-categories}, \autoref{sec:magnitude-homology}),
we showed that equipping case categories with \([0,1]\)-valued
hom-objects yields a principled bridge between symbolic case grammar and
statistical semantics. Categorical magnitude
(\autoref{sec:magnitude-homology}) provides a quantitative ``effective
size'\,' invariant for comparing case systems; magnitude homology
{[}-@leinster2021magnitude{]} refines that comparison when scalar
magnitude does not separate two systems.
\end{block}

\begin{block}{C5: Topos-Theoretic Transfer via Morita Equivalence}
\protect\phantomsection\label{c5-topos-theoretic-transfer-via-morita-equivalence}
Using Caramello's {[}-@caramello2016bridges; -@caramello2021five{]}
bridge technique and Phillips's {[}-@phillips2024lot{]} universality
result (\autoref{sec:topos-theory}), we articulate how
\textbf{topos-theoretic invariants} of a classifying topos can scaffold
inter-theoretic transfer: when two formalizations admit
Morita-equivalent classifying toposes (or an explicit bridge topos, as
sketched in \autoref{sec:topos-theory}), properties expressible as
shared invariants transfer without separate proof in each framework. The
schematic equivalence chain for typological, type-logical,
distributional, and enriched case theories is a \textbf{research
program}---not a single finished theorem covering all of
linguistics---but it aligns the four perspectives with Caramello's
methodology.
\end{block}

\begin{block}{C6: Diagrams as Cognitively Privileged Representations}
\protect\phantomsection\label{c6-diagrams-as-cognitively-privileged-representations}
The meta-contribution (Sections 1 and 7) is the argument---supported by
the cognitive science of diagrams {[}@larkin1987diagram;
@shimojima1996reasoning; @giardino2017diagrammatic;
@manders2008euclidean{]}---that commutative diagrams are not merely
convenient notation but cognitively privileged representations that
provide inferential advantages in case reasoning. When embedded in an
active inference framework {[}@namjoshi2026fundamentals{]} and the
CEREBRUM architecture {[}@friedman2024cerebrum{]}, these diagrams serve
as the structural core of generative models for total cognitive scenario
understanding.
\end{block}

\begin{block}{C7: Computational Verification and Test Suite
(\textbf{implemented and tested})}
\protect\phantomsection\label{c7-computational-verification-and-test-suite-implemented-and-tested}
The framework is computationally implemented and verified through
\textbf{1197} automated tests across \textbf{64} test files (sixty-four
files). 95.96\% line-and-branch coverage on \texttt{src/} (from
\texttt{coverage.json}) The CI/build gate requires \textbf{≥90\%}
coverage on \texttt{src/} (see \autoref{sec:diagrammatic-cognition},
DAIF in \autoref{sec:daif-results}, and
\texttt{src/generate\_manuscript\_metrics.py} for injection of these
values at build time). The \texttt{src/daif/} subpackage (7 modules, 25
symbols, 224 dedicated tests) provides distributional RL components
(push-forward, quantile TD, VMP, Bethe FE, ERP profiles, policy
selection, diagnostics). All tests use real mathematical
computations---no mocks---ensuring results reflect genuine behaviour of
the formal structures in \texttt{src/}. Raw manuscript sources contain
\texttt{\$\{...\}} placeholders resolved by the metrics pipeline.
\end{block}

\begin{block}{C8: Quantum Active Inference and Topological Semantic Flow
(\textbf{theoretical bridge})}
\protect\phantomsection\label{c8-quantum-active-inference-and-topological-semantic-flow-theoretical-bridge}
The framework's categorical string diagrams connect to literature on
topological quantum neural networks, the ZX-calculus, and
sheaf-theoretic quantum semantic communication
(\autoref{sec:quantum-active-inference}). The \texttt{src/quantum/}
module implements a POVM measurement model for case roles
(\textbf{implemented and tested}). The broader TQNN spin-networks, full
ZX circuit compilation, and hardware claims represent a theoretical
bridge rather than complete local execution in this repository.
\end{block}

\begin{block}{C9: Cognitive Security and Case-Theoretic AI Safety
(\textbf{specification and proxy implementation})}
\protect\phantomsection\label{c9-cognitive-security-and-case-theoretic-ai-safety-specification-and-proxy-implementation}
The framework provides a case-theoretic analysis of AI security
(\autoref{sec:cognitive-security}; see also compositional protocol
typing in \autoref{sec:ai-implications}). Prompt injection is analyzed
as structurally equivalent to illicit case-role re-assignment (type
violation in the interaction grammar; \textbf{implemented as
type-checking proxy in \texttt{src/security/} and
\texttt{src/diagrams/}}). This leads to the proposal of case-theoretic
firewalls and epistemic case security as a framework. The
\texttt{src/security/cognitive\_security.py} implements finite-category
validation and scoring (tested). Full production enforcement and
non-cartesian monoidal constraints remain open engineering targets, as
detailed in the limitations section of \autoref{sec:cognitive-security}.
\end{block}

\begin{block}{C10: Falsifiable Neurolinguistic Predictions}
\protect\phantomsection\label{c10-falsifiable-neurolinguistic-predictions}
The integration of enriched category theory with active inference
generates quantitative, testable predictions about neural processing of
case structure (\autoref{sec:cognitive-integration}). Specifically, the
amplitude of prediction-error ERP components (P600, N400) is predicted
to scale with the enriched hom-value (precision) of the violated
morphism, and garden-path reanalysis costs are predicted to correlate
with the change in categorical magnitude between the initial and revised
case diagrams. These hypotheses \textbf{extend} computational accounts
that decompose surprisal into components linked to N400- vs P600-like
signatures {[}@li2023decomposition; @li2024shallow{]}, by tying
amplitudes to enriched morphism weights and magnitude change in an
explicit case-diagram prior. They bridge abstract categorical formalism
and empirical neurolinguistics, making the framework falsifiable at the
single-trial level.
\end{block}

\begin{block}{C11: Categorical Communication Protocols for Multi-Agent
AI}
\protect\phantomsection\label{c11-categorical-communication-protocols-for-multi-agent-ai}
The synthesis of case categories with modern agent communication
standards (A2A, MCP, ACP, ANP) yields a principled design for
compositional, type-safe agent protocols
(\autoref{sec:ai-implications}). By assigning case roles (NOM, ACC, INS,
LOC, DAT) to interaction participants and enforcing relational type
constraints at protocol boundaries, the framework provides
interpretable, composable interaction schemas that go beyond the flat
JSON payloads of current standards, with topos-level invariants
transferable across Morita-equivalent formalizations when equivalence is
exhibited (\autoref{sec:topos-theory}). The DisCoCirc extension enables
discourse-level tracking of agent states across multi-turn dialogues,
with protocol correctness reducing to categorical type-checking.
\end{block}
\end{frame}

\begin{frame}[fragile,allowframebreaks]{Eight Open Directions}
\protect\phantomsection\label{sec:future-directions}
The following eight directions (\textbf{F1}--\textbf{F8}) identify the
most tractable and impactful extensions of this framework:

\begin{block}{F1: Computational Experiments with DisCoCirc and lambeq}
\protect\phantomsection\label{f1-computational-experiments-with-discocirc-and-lambeq}
The DisCoCirc framework {[}@defelice2022discocirc{]} offers a natural
platform for testing our case-theoretic predictions computationally. The
release of \textbf{lambeq Gen II} {[}@lambeq2025genii{]} with full
DisCoCirc support makes this direction immediately tractable:
discourse-level case role tracking---including the dynamic role
reversals of \autoref{fig:three-sentence-discourse}---can now be
compiled into parameterized quantum circuits (PQCs) and trained
end-to-end. Krawchuk et al. {[}@krawchuk2025paramqnlp{]} demonstrate
efficient generation of PQCs from large-scale texts (up to 6,410 words)
with competitive performance on sentiment classification;
{[}@letcher2024tight{]} and {[}@rad2024trainability{]} provide gradient
bounds and reduced-domain initialization techniques that mitigate barren
plateaus, making discourse-level circuit training practically feasible.
Concrete experiments could include:

\begin{itemize}
\tightlist
\item
  Implementing case-marked DisCoCat models in \textbf{lambeq Gen II} and
  evaluating them on semantic role labeling (SRL) tasks, leveraging the
  discourse-level wiring to track case roles across sentence boundaries
\item
  Comparing the accuracy of alignment-specific meaning functors
  (accusative vs.~ergative) on typologically diverse corpora
\item
  Measuring the categorical magnitude of empirically derived enriched
  case categories and correlating it with typological complexity
  measures
\item
  Using the \texttt{complexity\_metrics} module to quantify derivational
  complexity across sentence types and correlating syntactic complexity
  scores with human processing difficulty (reading times, surprisal)
\item
  Training lambeq Gen II circuits on case role reversal discourses using
  reduced-domain parameter initialization {[}@rad2024trainability{]} to
  avoid local minima
\end{itemize}
\end{block}

\begin{block}{F2: Topos-Theoretic Grammatical Induction from Corpora}
\protect\phantomsection\label{f2-topos-theoretic-grammatical-induction-from-corpora}
Caramello's {[}-@caramello2023syntactic{]} syntactic learning technique
could be applied to induce case-theoretic axioms from annotated corpora.
The program would:

\begin{enumerate}
\tightlist
\item
  Extract case-labeled dependency structures from a Universal
  Dependencies treebank
\item
  Construct the classifying topos of the implicit case theory
\item
  Read off the alignment type, morphism structure, and enriched weights
  from the topos-theoretic axioms
\item
  Compare the induced theory against typological descriptions
\end{enumerate}
\end{block}

\begin{block}{F3: Quantum Case Categories on Near-Term Hardware}
\protect\phantomsection\label{f3-quantum-case-categories-on-near-term-hardware}
The QNLP connection (\autoref{sec:categorical-semantics}) opens the
possibility of implementing case categories directly as quantum
circuits. Case roles would correspond to quantum registers, grammatical
relations to parameterized gates, and alignment functors to circuit
transformations. This would provide a genuinely new computational
paradigm for grammatical inference---exploiting quantum parallelism to
explore the space of case assignments simultaneously. Two recent results
make this direction practically tractable: (1) Rad et al.'s
{[}-@rad2024trainability{]} reduced-domain parameter initialization
yields polynomial gradient decay, suppressing the barren plateau problem
for circuits of linguistic depth; and (2) Letcher et al.'s
{[}-@letcher2024tight{]} assumption-free gradient bounds rule out
vanishing gradients for circuits with local observables, which includes
the case-role measurement POVMs of \autoref{eq:eq-8-1}. Together these
results suggest that near-term quantum hardware can support case
category training without exponential gradient overhead.
\end{block}

\begin{block}{F4: Neural Predictive Processing and Electrophysiological
Predictions}
\protect\phantomsection\label{f4-neural-predictive-processing-and-electrophysiological-predictions}
The predictive processing account of \autoref{sec:cognitive-integration}
generates testable neuroscientific predictions:

\begin{itemize}
\tightlist
\item
  Case-marking violations should elicit prediction-error responses
  (P600/N400) proportional to the enriched weight of the violated
  morphism
\item
  Typologically unusual case patterns should require more precision
  updating than expected patterns
\item
  Diagrammatic representations of case structure should be decodable
  from neural activity during sentence comprehension
\end{itemize}
\end{block}

\begin{block}{F5: Cross-Modal Case Structure in Embodied Cognition}
\protect\phantomsection\label{f5-cross-modal-case-structure-in-embodied-cognition}
The situation semantics connection (\autoref{sec:cognitive-integration})
suggests extending case categories beyond language to multi-modal
perception. Visual scene understanding also requires assigning
relational roles---who is acting on what, where things are located, what
instruments are being used. An active inference agent should maintain a
unified case diagram that integrates linguistic, visual, and motor
information, using the same categorical structure for all modalities.
This would provide a formal account of how language grounds in
perception and action---a key challenge for embodied AI.
\end{block}

\begin{block}{F6: Enriched Category Learning from Distributional Data}
\protect\phantomsection\label{f6-enriched-category-learning-from-distributional-data}
Bradley's {[}-@bradley2024ipam; -@bradley2025tea{]} program of treating
language itself as an enriched category suggests a learning algorithm:
estimate the enriched hom-values from corpus data and then extract the
categorical structure that best explains the observed distributional
patterns. Applied to case, this would yield \emph{empirically grounded}
case categories whose objects (roles) and morphisms (relations) emerge
from data rather than being stipulated a priori.
\end{block}

\begin{block}{F7: Extending Distributional Active Inference for
Linguistic Agents}
\protect\phantomsection\label{f7-extending-distributional-active-inference-for-linguistic-agents}
The Distributional Active Inference (DAIF) framework of Akgül et al.
{[}-@akgul2026distributional{]} has been computationally implemented in
this paper's \texttt{src/daif/} subpackage, integrating active inference
into distributional reinforcement learning
{[}@bellemare2017distributional{]} via push-forward measures on
representation paths (\autoref{sec:daif-results}). The current
implementation models case assignment distributionally: agents maintain
\emph{distributions over case diagrams}, sharpening beliefs through
variational message passing as linguistic evidence accumulates. Key open
extensions include:

\begin{itemize}
\tightlist
\item
  Training distributional case-assignment circuits end-to-end in
  \textbf{lambeq Gen II} using quantile regression losses, enabling
  gradient-based learning of the enriched weight matrix from annotated
  corpora
\item
  Extending the DAIF policy selection (\texttt{G\_policy},
  \texttt{softmax\_policy\_selection}) to multi-turn dialogue management
  with DisCoCirc entity persistence---tracking agent state distributions
  across sentence boundaries
\item
  Cross-lingual transfer of Bethe free energy convergence profiles:
  testing whether the free energy minima differ systematically between
  nominative-accusative and ergative-absolutive languages,
  operationalizing the typological complexity predictions of
  \autoref{sec:enriched-categories}
\item
  Integrating IQN risk distortion modes (optimistic/pessimistic/CVaR)
  into the CEREBRUM architecture to model individual differences in
  syntactic risk tolerance
\end{itemize}
\end{block}

\begin{block}{F8: Synthesizing Biolinguistic Syntax with Neuropragmatic
Inference via the ROSE Model}
\protect\phantomsection\label{f8-synthesizing-biolinguistic-syntax-with-neuropragmatic-inference-via-the-rose-model}
A critical open direction is fully computationalizing the handoff
between the rigid algebraic geometry of syntax and the highly
associative probabilistic inference of discourse. Murphy's \textbf{ROSE}
(Representation, Operation, Structure, Encoding) model
{[}@murphy2023rose{]} suggests that the brain achieves this via
cross-frequency phase-amplitude coupling (PAC), establishing a
``mesoscopic protectorate'' for formal operations. Future extensions
should:

\begin{itemize}
\tightlist
\item
  Model PAC directly within CEREBRUM by assigning distinct temporal
  decay rates to syntactic operators vs.~pragmatic context vectors.
\item
  Simulate the export of case-marked commutative diagrams from a
  simulated ``core language network'' to a simulated Default Mode
  Network using Gutiérrez Cisneros et al.'s
  {[}-@gutierrez2026neuropragmatics{]} framework for speech-act
  evaluation.
\item
  Test whether the diagrammatic free-ride inferences
  (\autoref{fig:case-minimal}) predicted by our categorical model map
  onto the theta-gamma cross-frequency signatures observed during
  pragmatic garden-path recovery.
\end{itemize}
\end{block}
\end{frame}

\begin{frame}[allowframebreaks]{Five Takeaways: One Argumentative Line
per Formal Pillar}
\protect\phantomsection\label{five-takeaways-one-argumentative-line-per-formal-pillar}
For the reader who absorbs nothing else, the argumentative line reduces
to these five points --- one per pillar of the formal architecture:

\begin{enumerate}
\item
  \textbf{Case is category-theoretic, not morphological.} The universal
  fact of linguistic case is that every language wires up \emph{who did
  what to whom}; the formal content of that wiring is precisely the data
  of a category with case roles as objects and grammatical relations as
  morphisms. Alignment typologies (nominative, ergative, tripartite,
  active-stative, fluid-S) are structure-preserving functors between
  case categories --- a typological taxonomy becomes a proof-by-functor.
\item
  \textbf{String diagrams make the grammar executable.} Lambek pregroup
  types, compact closure, and the snake equation compile each sentence
  into a DisCoPy string diagram whose normal form is the sentence type
  \(s\); each reduction is a proof of well-typedness, and the reduction
  \emph{is} the diagram. The three-sentence DisCoCirc example ships a
  working discourse with per-entity role-history ribbons showing Alice
  traversing NOM → ACC → NOM.
\item
  \textbf{Enriched weights ground graded intuitions.} Moving from
  \(\mathbf{Bool}\)- to \([0,1]\)-enrichment replaces binary
  admissibility with graded distributional proximity, makes categorical
  magnitude \(|\mathcal{C}| \approx 2.5\) a computable summary of
  effective role distinctions, and gives the enriched weight
  \(w_f = \mathcal{C}(A, B)\) used throughout
  \autoref{sec:cognitive-integration} / \autoref{sec:daif-results} as
  the precision on every prediction error.
\item
  \textbf{DAIF ports the whole scaffolding into active inference.}
  Variational message passing, the Bethe free energy, and a four-term
  expected free energy \(G(\pi)\) turn distributional case assignment
  into a first-principles inference loop. The N400 and P600 ERP
  amplitudes fall out of a first-order expansion
  \(\Delta F \approx -\Delta \mathbb{E}[Z] + \tfrac{1}{2}\Delta\Lambda\sigma_Z^{2}\)
  --- derivations, not empirical ansätze.
\item
  \textbf{The same formal layers give AI-safety researchers three
  concrete handles.} Type discipline on messages, decidable
  admissibility of multi-turn sequences, and graded-confidence
  attenuation (\autoref{sec:ai-implications}) give prompt injection a
  type-theoretic reformulation (\autoref{sec:cognitive-security}) --- an
  engineering specification for agent protocols, not a guarantee on
  untyped present-day LLM APIs.
\end{enumerate}
\end{frame}

\begin{frame}[fragile,allowframebreaks]{What the Paper Does \emph{Not}
Claim: Consolidated Limitations}
\protect\phantomsection\label{what-the-paper-does-not-claim-consolidated-limitations}
Four limitations recur across the framework and are recorded here
explicitly rather than left implicit in their sections of origin:

\begin{itemize}
\tightlist
\item
  \textbf{The mean-field approximation in
  \texttt{push\_forward\_return}} maintains one belief-weighted return
  distribution instead of per-state distributions \(Z(s)\). By the
  mean-field bound recorded in \autoref{sec:daif-limitations} the
  approximation error is at most \(\gamma \cdot R_{\max} \cdot H[q]\) in
  the \(W_1\) metric --- tight for sharp posteriors, linearly degrading
  as entropy grows.
\item
  \textbf{The enriched-categorical unification is a conjecture.}
  Distributional semantics, distributional RL, and active-inference
  posteriors all instantiate \([0,1]\)-enriched structures, but a strict
  categorical proof that the three share a common enriching monoidal
  base remains open (\autoref{sec:daif-limitations}).
\item
  \textbf{Empirical validation is narrow.} Case-assignment
  demonstrations use a single German sentence. Cross-linguistic and
  cross-register generalisation is left to future work; the hooks in
  \texttt{make\_daif\_belief\_trajectory\_data()} make adding corpora a
  one-function change.
\item
  \textbf{ERP amplitudes are not calibrated to μV.} The DAIF predictions
  are on a dimensionless return/log-probability scale; converting to μV
  would require a per-subject scaling constant fit to empirical ERP
  data. The qualitative ordering and graded precision response are
  predictions the current framework \emph{does} make.
\end{itemize}
\end{frame}

\begin{frame}[fragile,allowframebreaks]{Anticipated Objections and
Responses}
\protect\phantomsection\label{sec:reviewer-objections}
Four lines of objection stand out; we address each honestly rather than
waiting for a reviewer.

\textbf{Objection 1 --- ``The mean-field approximation throws away the
whole point of distributional RL.''} \emph{Response.} The
belief-weighted collapse gives \(\mathcal{O}(n)\) memory in place of
\(\mathcal{O}(n \cdot N_{\text{atoms}})\) and is exact in the
sharp-posterior limit. The mean-field bound recorded in
\autoref{sec:daif-limitations} quantifies the degradation as \(H[q]\)
grows. The alternative --- maintaining per-state return distributions
throughout a linguistic parse --- would multiply the runtime of
\texttt{distributional\_case\_assignment()} by a factor of \(n\) (the
role count) for a gain that vanishes as the posterior sharpens after the
first few words of a sentence.

\textbf{Objection 2 --- ``The enriched-categorical unification is only
stated, never proven.''} \emph{Response.} Conceded, and flagged in
\autoref{sec:daif-limitations} and in the What's New block of
\autoref{sec:whats-new}. Distributional semantics, distributional RL,
and active-inference posteriors all \emph{instantiate}
\([0,1]\)-enriched structure; whether they share a common enriching
monoidal base in the strict sense of Kelly enriched-category theory is
an open question. The repository's tests verify the enriched axioms hold
in each instance; they do not construct the common base functor, and the
manuscript nowhere uses the strict unification as a load-bearing step in
any downstream argument.

\textbf{Objection 3 --- ``Cross-linguistic evidence is a single German
sentence.''} \emph{Response.} The empirical scope reflects the
publication's focus on \emph{specification} rather than \emph{corpus
study}. The \texttt{Discourse.role\_reversal("Alice",\ "Bob")} and
\texttt{make\_daif\_belief\_trajectory\_data()} interfaces accept
arbitrary lexical heads and observation sequences; extending to Basque,
Dyirbal, Russian, Serbian/BCS, or any other language is a one-function
change. \autoref{fig:multilingual-isomorphism} (in
\autoref{sec:categorial-grammar}) already shows a pregroup multilingual
isomorphism across English / Latin / Japanese, and the Slavic discussion
in \autoref{sec:case-systems} extends the same type calculus to overtly
case-marked Russian and Serbian noun phrases. Slavic-language ERP
datasets --- Bornkessel-Schlesewsky and colleagues' eADM-aligned Russian
case-violation paradigms in particular {[}@bornkessel2006extended{]} ---
would provide a near-term empirical test for the precision-weighted P600
prediction in \textbf{C10} / \textbf{F4}. The three-sentence Alice/Bob
discourse is a proof-of-concept, not a typological claim.

\textbf{Objection 4 --- ``The ROSE phase--amplitude coupling gap means
you cannot predict ERP latencies.''} \emph{Response.} Conceded and
flagged explicitly in \autoref{sec:daif-limitations}. The current
implementation keeps N400 and P600 \emph{peak latencies} as fixed
Gaussian peaks at 380 ms and 600 ms respectively (see
\texttt{DEFAULT\_N400\_PEAK\_MS} / \texttt{DEFAULT\_P600\_PEAK\_MS} in
\texttt{src/daif/prediction.py}). A principled latency prediction would
require a cross-frequency-coupling delay parameter inside the CEREBRUM
layer; we flag this as the cleanest next research step.
\end{frame}

\begin{frame}[allowframebreaks]{What to Read Next, by Reader Profile}
\protect\phantomsection\label{what-to-read-next-by-reader-profile}
Readers who entered this paper from different fields will want different
onward paths:

\begin{itemize}
\tightlist
\item
  \textbf{Linguists / typologists}: skip ahead to
  \autoref{sec:case-systems}--\autoref{sec:case-categories} for the
  categorical formalisation of case and the five alignment functors;
  then \autoref{sec:compact-closure-complexity} for the complexity
  metrics that make cross-linguistic comparison quantitative.
\item
  \textbf{Machine-learning researchers}: \autoref{sec:daif-results}
  first for the Distributional Active Inference extension of the
  Bellemare-Dabney-Munos programme to linguistic case; then
  \autoref{sec:daif-erp} for how DAIF yields falsifiable ERP
  predictions.
\item
  \textbf{Cognitive / computational neuroscientists}:
  \autoref{sec:cognitive-integration} for the active-inference framing;
  \autoref{sec:diagrammatic-cognition} for the three falsifiable ERP
  predictions; \autoref{sec:daif-limitations} for the PAC-latency gap
  that is the cleanest experimental entry point.
\item
  \textbf{QNLP / quantum-semantics engineers}:
  \autoref{sec:quantum-active-inference} and
  \autoref{sec:quantum-semantics} for the POVM-based case assignment and
  the honest separation of implemented POVM machinery from
  literature-bridge TQNN / ZX / lambeq claims.
\item
  \textbf{AI-safety / agent-protocol practitioners}:
  \autoref{sec:ai-implications} for the three concrete handles (type
  discipline, decidable admissibility, graded confidence) and
  \autoref{sec:cognitive-security} for prompt injection as a type
  violation --- engineering specification rather than automatic
  guarantee on today's APIs.
\end{itemize}
\end{frame}

\begin{frame}[allowframebreaks]{Case Categories Are the Geometry of
Meaning: A Unifying Coda}
\protect\phantomsection\label{case-categories-are-the-geometry-of-meaning-a-unifying-coda}
The commutative diagram is the central motif of this review---both as a
mathematical tool and as a cognitive instrument. The analysis has
demonstrated that the same diagrammatic language that makes category
theory effective for formalizing case systems also makes it effective
for \emph{thinking about} case systems: the spatial structure of a
commutative diagram encodes relational information in a format that
supports rapid search, pattern recognition, and free-ride inference.

This convergence of mathematical utility and cognitive efficacy is not
accidental. If the active inference framework is correct, then the brain
operates by constructing and updating generative models of the world's
relational structure. Category theory provides the structural algebra
for these generative models; commutative diagrams supply their natural
topology; and case categories instantiate the precise relational
vocabulary that cognitive systems deploy to organize experience into
coherent narratives of \emph{who does what to whom}. The distributional
revolution in both semantics and reinforcement learning---from Firth's
{[}-@firth1957papers{]} co-occurrence statistics through transformer
attention weights to Distributional Active Inference
{[}@akgul2026distributional{]}---confirms that meaning is not an atomic
property of symbols but emerges from the relational structure of
contexts, a principle that enriched category theory captures with
mathematical precision.

Crucially, \textbf{this synthesis clarifies a mathematical angle on
alignment for relational agent cognition.} Moving from flat token
streams toward \textbf{explicitly typed} enriched categorical
interaction grammars lets one state \textbf{when} prompt injection
corresponds to a \textbf{type error relative to a fixed protocol}---the
conditional analysis in \autoref{sec:cognitive-security}. Default LLM
stacks do not yet enforce such protocols end-to-end; the payoff is
sharper \textbf{specification} (what boundary checks would guarantee)
and a research agenda for non-cartesian wiring, not a blanket claim that
injection is already impossible in production systems.

Ultimately, the mathematics of case alignment presents a highly
structured formal geometry of meaning---the relational algebra through
which cognitive agents navigate and render intelligible the structured
world of events, participants, and relations. The convergence of formal
semantics, distributional semantics, and active inference within a
single categorical framework suggests that commutative diagrams offer a
natural formalism in which relational cognition can be modeled.
\end{frame}

\end{document}
